\chapter{อุณหพลศาสตร์ออกซิเจนและการควบคุมเครื่องตีน้ำประหยัดพลังงาน}
\label{chap:ch06}

\section*{แผนบริหารการสอนประจำบทที่ 6}
\addcontentsline{toc}{section}{แผนบริหารการสอนประจำบทที่ 6}

\begin{planbox}{หัวข้อเนื้อหาประจำบท}
\begin{enumerate}[topsep=2pt, itemsep=1pt, partopsep=0pt, parsep=0pt, leftmargin=1.5em]
    \item สรีรวิทยาการหายใจของกุ้งขาวแวนนาไมและความสำคัญของออกซิเจนละลายน้ำ
    \item แบบจำลองอุณหพลศาสตร์ความเข้มข้นอิ่มตัวของออกซิเจนของเบนสัน-เคราส์
    \item จลนพลศาสตร์การถ่ายเทมวลสารของก๊าซและอัตราการถ่ายเทออกซิเจนมาตรฐาน
    \item กฎความคล้ายคลึงของเครื่องกลหมุนและการประหยัดพลังงานไฟฟ้าด้วยอินเวอร์เตอร์ VFD
    \item อัลกอริทึมการควบคุมเครื่องตีน้ำแบบลีด-แล็กและการสลับการทำงาน 4 ลำดับขั้น
\end{enumerate}
\end{planbox}

\begin{planbox}{วัตถุประสงค์เชิงพฤติกรรม}
\begin{enumerate}[topsep=2pt, itemsep=1pt, partopsep=0pt, parsep=0pt, leftmargin=1.5em]
    \item อธิบายผลกระทบของอุณหภูมิและความเค็มต่อค่าความเข้มข้นอิ่มตัวของออกซิเจนในน้ำได้
    \item คำนวณค่าออกซิเจนอิ่มตัวตามสมการเบนสัน-เคราส์ได้อย่างถูกต้อง
    \item ประเมินอัตราการถ่ายเทออกซิเจน (OTR) และค่าสัมประสิทธิ์การถ่ายเทมวลสาร ($K_L a$) ได้
    \item ออกแบบอัลกอริทึมควบคุมเครื่องตีน้ำแบบวงปิดเพื่อประหยัดพลังงานไฟฟ้าร้อยละ 40 ได้
\end{enumerate}
\end{planbox}

\newpage

\begin{planbox}{วิธีสอนและกิจกรรมการเรียนการสอน}
\begin{enumerate}[topsep=2pt, itemsep=1pt, partopsep=0pt, parsep=0pt, leftmargin=1.5em]
    \item บรรยายเชิงทฤษฎีอุณหพลศาสตร์การละลายของก๊าซและชลศาสตร์การกวนน้ำ
    \item สาธิตการทดลองวัดค่า DO ตลอด 24 ชั่วโมงในบ่อเลี้ยงกุ้งจำลอง
    \item ปฏิบัติการเขียนโค้ด Python จำลองสมการเบนสัน-เคราส์และคำนวณกราฟประหยัดพลังงาน
    \item ออกแบบและทดสอบโปรแกรมควบคุมการสลับมอเตอร์ Lead-Lag\index{การควบคุมแบบลีด-แล็ก (Lead-Lag Control)}\index{Lead-Lag Control} บนระบบจำลอง
\end{enumerate}
\end{planbox}

\begin{planbox}{สื่อการเรียนการสอน}
\begin{enumerate}[topsep=2pt, itemsep=1pt, partopsep=0pt, parsep=0pt, leftmargin=1.5em]
    \item เอกสารคำสอนบทที่ 6 สไลด์บรรยาย และตารางความเข้มข้นอิ่มตัวของออกซิเจน
    \item ชุดจำลองถังน้ำทดลองเติมอากาศพร้อมหัววัด Optical DO และหัววัดความเค็ม
    \item อินเวอร์เตอร์ขับมอเตอร์ VFD และมอเตอร์ใบพัดตีน้ำจำลอง
    \item สคริปต์แบบจำลองคอมพิวเตอร์คำนวณการใช้พลังงานไฟฟ้า
\end{enumerate}
\end{planbox}

\begin{planbox}{การวัดและประเมินผล}
\begin{enumerate}[topsep=2pt, itemsep=1pt, partopsep=0pt, parsep=0pt, leftmargin=1.5em]
    \item การทดสอบข้อเขียนเรื่องสมการเบนสัน-เคราส์และการถ่ายเทออกซิเจน
    \item การประเมินความถูกต้องของแบบฝึกหัดการคำนวณกำลังไฟฟ้าตามกฎความคล้ายคลึง
    \item การประเมินความสมบูรณ์ของโค้ดระบบควบคุมเครื่องตีน้ำแบบ 4 ลำดับขั้น
\end{enumerate}
\end{planbox}

\clearpage

\section{สรีรวิทยาการหายใจและระดับออกซิเจนวิกฤติในบ่อกุ้ง}

ในการเพาะเลี้ยงกุ้งขาวแวนนาไมหนาแน่นสูง (Intensive Aquaculture) ออกซิเจนละลายน้ำ (DO) เป็นปัจจัยจำกัดการเจริญเติบโตที่สำคัญที่สุด ระดับออกซิเจนที่เหมาะสมสำหรับการเจริญเติบโตของกุ้งขาวต้องมากกว่า \textbf{5.0 mg/L} 
\begin{itemize}[leftmargin=1.5em]
    \item หากค่า DO ลดต่ำลงมาอยู่ในช่วง $3.0 - 4.0\text{ mg/L}$ กุ้งจะเริ่มหยุดกินอาหาร อัตราการเปลี่ยนอาหารเป็นเนื้อ (FCR) จะแย่ลง และภูมิคุ้มกันโรคจะดรอปลงอย่างรวดเร็ว
    \item หากค่า DO ลดต่ำกว่า \textbf{2.5 mg/L} กุ้งจะเกิดอาการขาดออกซิเจนรุนแรง ว่ายลอยหัวตามผิวน้ำ และอาจเกิดการตายยกบ่อได้ภายในเวลาไม่กี่ชั่วโมง
\end{itemize}

\section{แบบจำลองอุณหพลศาสตร์ออกซิเจนอิ่มตัวของเบนสัน-เคราส์}

ความสามารถในการละลายของก๊าซออกซิเจนในน้ำขึ้นอยู่กับอุณหภูมิ ความดันบรรยากาศ และความเค็ม \textbf{แบบจำลองเบนสัน-เคราส์ (Benson and Krause Model)\index{แบบจำลองเบนสัน-เคราส์ (Benson-Krause)}\index{Benson-Krause Model}} เป็นสมการมาตรฐานสากลในการคำนวณความเข้มข้นอิ่มตัวของออกซิเจน ($C^*_s$) ในน้ำกร่อยและน้ำเค็ม

\begin{equation}
\ln C^*_s = A_1 + \frac{A_2}{T} + \frac{A_3}{T^2} + \frac{A_4}{T^3} - S \left( B_1 + \frac{B_2}{T} + \frac{B_3}{T^2} \right)
\label{eq:benson_krause}
\end{equation}

โดยที่ $T$ แทนอุณหภูมิสัมบูรณ์ในหน่วยเคลวิน ($T = t\,^\circ\text{C} + 273.15$), $S$ แทนค่าความเค็มในหน่วยส่วนในพันส่วน (ppt หรือ g/kg), และค่าคงที่การละลายกำหนดโดย $A_1 = -139.34411$, $A_2 = 1.575701 \times 10^5$, $A_3 = -6.642308 \times 10^7$, $A_4 = 1.243800 \times 10^{10}$, $B_1 = -0.017674$, $B_2 = 10.754$, และ $B_3 = -2140.7$

\begin{table}[H]
\centering
\caption{ความเข้มข้นอิ่มตัวของออกซิเจนในน้ำ ($C^*_s$, mg/L) ที่ระดับความเค็มและอุณหภูมิต่างๆ}
\label{tab:do_saturation}
\begin{tabularx}{\textwidth}{cYYYY}
    \toprule
    \textbf{อุณหภูมิ ($^\circ\text{C}$)} & \textbf{น้ำจืด (0 ppt)} & \textbf{น้ำกร่อย (15 ppt)} & \textbf{น้ำทะเล (25 ppt)} & \textbf{น้ำทะเลเข้มข้น (35 ppt)} \\
    \midrule
    26 & 8.11 & 7.42 & 6.99 & 6.58 \\
    28 & 7.83 & 7.17 & 6.76 & 6.37 \\
    30 & 7.56 & 6.93 & 6.54 & 6.17 \\
    32 & 7.31 & 6.71 & 6.33 & 5.98 \\
    34 & 7.07 & 6.49 & 6.13 & 5.79 \\
    \bottomrule
\end{tabularx}
\end{table}

\section{จลนพลศาสตร์การถ่ายเทมวลสารและอัตราการเติมออกซิเจน}

อัตราการเพิ่มขึ้นของออกซิเจนละลายน้ำในสภาวะที่ไม่มีสิ่งมีชีวิตถูกอธิบายโดยกฎการแพร่สองฟิล์ม (Two-Film Theory) ดังสมการ

\begin{equation}
\frac{dC}{dt} = K_L a \left( C^*_s - C \right)
\label{eq:oxygen_transfer}
\end{equation}

โดยที่ $K_L a$ แทนค่าสัมประสิทธิ์การถ่ายเทมวลสารโดยรวม ($\text{hr}^{-1}$), $C$ แทนค่าออกซิเจนละลายที่วัดได้จริง, และ $C^*_s - C$ แทนแรงขับเคลื่อนของการละลาย (Concentration Deficit)

เมื่อเปิดเครื่องตีน้ำในบ่อเลี้ยงกุ้งที่มีออกซิเจนสูงอยู่แล้ว แรงขับเคลื่อน $C^*_s - C$ จะมีค่าน้อยมาก ทำให้การเติมออกซิเจนแทบไม่เกิดขึ้นและเป็นการสิ้นเปลืองพลังงานไฟฟ้าไปโดยเปล่าประโยชน์

\section{การควบคุมเครื่องตีน้ำประหยัดพลังงานด้วย VFD และระบบ Lead-Lag}

การใช้พลังงานไฟฟ้าของเครื่องตีน้ำใบพัดเป็นไปตาม \textbf{กฎความคล้ายคลึงของเครื่องกลหมุน (Affinity Laws)\index{กฎความคล้ายคลึงของเครื่องกลหมุน (Affinity Laws)}\index{Affinity Laws}} ซึ่งระบุว่ากำลังไฟฟ้าที่ใช้ ($P$) จะแปรผันตรงกับ \textbf{ความเร็วรอบยกกำลังสาม}

\begin{equation}
\frac{P_2}{P_1} = \left( \frac{N_2}{N_1} \right)^3
\label{eq:affinity_laws}
\end{equation}

เมื่อปรับลดความเร็วรอบของเครื่องตีน้ำลงเหลือร้อยละ 80 (ความถี่ $40\text{ Hz}$ จาก $50\text{ Hz}$) กำลังไฟฟ้าที่มอเตอร์ใช้จะลดลงเหลือเพียง $(0.8)^3 \approx 0.512$ หรือประหยัดพลังงานไฟฟ้าลงได้ถึง \textbf{ร้อยละ 48.8} ในขณะที่ยังคงสามารถสร้างกระแสน้ำหมุนวนรวมตะกอนเลนก้นบ่อได้เป็นอย่างดี

\begin{theorybox}[ระบบควบคุมการทำงาน 4 ลำดับขั้นแบบสลับมอเตอร์ Lead-Lag]
เพื่อรักษาอายุการใช้งานของมอเตอร์และประหยัดพลังงานสูงสุด ระบบจะแบ่งการทำงานเป็น 4 สเต็ป:
\begin{itemize}[leftmargin=1.5em]
    \item \textbf{สเต็ป 1 (DO $> 6.5\text{ mg/L}$):} ช่วงบ่ายแดดจัด สาหร่ายสังเคราะห์แสงได้ออกซิเจนสูง เปิดเครื่องตีน้ำตัวหลักที่ความถี่ $35\text{ Hz}$ เพื่อรวมตะกอนเลน
    \item \textbf{สเต็ป 2 ($5.0 < \text{DO} \le 6.5\text{ mg/L}$):} ช่วงหัวค่ำ เปิดเครื่องตีน้ำ 2 ตัวที่ความถี่ $42\text{ Hz}$
    \item \textbf{สเต็ป 3 ($4.0 < \text{DO} \le 5.0\text{ mg/L}$):} ช่วงดึก เปิดเครื่องตีน้ำ 3 ตัวที่ความถี่พิกัด $50\text{ Hz}$
    \item \textbf{สเต็ป 4 ($\text{DO} \le 4.0\text{ mg/L}$):} สภาวะวิกฤติช่วงเช้ามืด เปิดเครื่องตีน้ำครบทุกตัวเต็มพิกัด $50\text{ Hz}$ พร้อมส่งสัญญาณเตือนภัยฉุกเฉินผ่าน LINE Notify
    \item มอเตอร์จะถูกสลับบทบาทระหว่างตัวหลัก (Lead) และตัวรอง (Lag) ทุก 24 ชั่วโมง เพื่อให้ชั่วโมงการทำงานเฉลี่ยสะสมของทุกตัวเท่ากัน
\end{itemize}
\end{theorybox}

\begin{pythonbox}[อัลกอริทึมจำลองการควบคุมเครื่องตีน้ำ 4 สเต็ปด้วย Python]
def multistage_aerator_controller(do_mg_l: float, motor_hours: list):
    # เลือกลำดับมอเตอร์ที่มีชั่วโมงทำงานน้อยที่สุดขึ้นมาเป็น Lead
    sorted_motors = sorted(range(len(motor_hours)), key=lambda i: motor_hours[i])
    
    if do_mg_l > 6.5:
        # สเต็ป 1 รัน 1 ตัว ประหยัดพลังงาน
        active_plan = [(sorted_motors[0], 35.0)]
    elif do_mg_l > 5.0:
        # สเต็ป 2 รัน 2 ตัว
        active_plan = [(sorted_motors[0], 42.0), (sorted_motors[1], 42.0)]
    elif do_mg_l > 4.0:
        # สเต็ป 3 รัน 3 ตัว
        active_plan = [(sorted_motors[0], 50.0), (sorted_motors[1], 50.0), (sorted_motors[2], 50.0)]
    else:
        # สเต็ป 4 วิกฤติ รันเต็มพิกัดทุกตัว
        active_plan = [(m_id, 50.0) for m_id in sorted_motors]
        
    return active_plan
\end{pythonbox}

\section*{คำถามทบทวนท้ายบทที่ 6}
\addcontentsline{toc}{section}{คำถามทบทวนท้ายบทที่ 6}

\begin{enumerate}
    \item จงอธิบายว่าเหตุใดความเข้มข้นอิ่มตัวของออกซิเจนละลายน้ำ ($C^*_s$) จึงลดลงเมื่ออุณหภูมิและความเค็มของน้ำเพิ่มสูงขึ้น
    \item จากสมการการถ่ายเทมวลสาร จงวิเคราะห์ว่าเหตุใดการเปิดเครื่องตีน้ำในเวลากลางวันที่มีแดดจัดจึงได้ประสิทธิภาพการเติมออกซิเจนต่ำกว่าเวลากลางคืน
    \item ตามกฎความคล้ายคลึงของเครื่องกลหมุน หากลดความถี่ของอินเวอร์เตอร์ VFD ขับมอเตอร์ตีน้ำจาก 50 Hz ลงเหลือ 40 Hz จะประหยัดพลังงานไฟฟ้าได้ร้อยละเท่าใด
    \item จงอธิบายหลักการทำงานของระบบสลับมอเตอร์ Lead-Lag ในการยืดอายุการใช้งานของมอเตอร์เครื่องตีน้ำในบ่อเลี้ยงสัตว์น้ำ
\end{enumerate}
\clearpage
